\section{Structured search and the granularity of a thought}
\label{sec:C09:search}
Tree of Thoughts treats intermediate text units as search states and uses generation, evaluation, and traversal rules to explore alternatives. Its Game of 24 experiment reports 74 percent success for the evaluated method versus 4 percent for its chain-of-thought comparator~\cite{yao2023tot}. Those author-reported results establish a task-specific system comparison with additional search; they do not isolate a universal architectural capability or equalize every inference cost. The inspected repository's \path{src/tot/methods/bfs.py} exposes \path{get_proposals}, \path{get_values}, \path{get_votes}, and \texttt{solve}~\cite{rewrite_tot_repo}. This makes proposal, evaluation, and search policy identifiable implementation components. In this snapshot, \path{get_value} caches evaluations by the rendered value prompt, while \path{get_values} assigns zero to repeated candidate strings within one evaluation call. These are not independent repeated checks. The sampled-selection branch in \path{solve} normalizes scores by their sum, so its mathematical domain requires nonnegative weights with positive total; an all-zero set needs an explicit fallback in an integration. Record candidate entries, unique artifacts, fresh evaluator calls and cache hits separately. Otherwise repeated or cached work can be mistaken for additional search evidence.

Granularity determines branching and verifiability. Token-level search has a large shallow action space and weak semantic boundaries. A complete-program proposal permits executable testing but can waste work when one early assumption is wrong. Intermediate structured states can expose constraints, but the state representation must preserve all information needed to continue. Merging two reasoning nodes because their summaries look similar can destroy an important distinction. Search-state equivalence needs a task-specific invariant, just as cache reuse requires a computational invariant in Chapter~\ref{ch:C04}.

With branching factor $b$ and depth $d$, an unpruned tree has $(b^{d+1}-1)/(b-1)$ nodes for $b\neq1$. Beam search with width $w$ can reduce retained nodes while pruning a correct path based on an imperfect early score. Depth-first traversal reduces frontier memory but can spend a long budget on one unproductive branch. Search order therefore trades coverage, peak storage, latency, and susceptibility to evaluator error. Report the actual explored graph and pruning decisions. A nominal beam width alone does not identify how much model computation was used when proposals and evaluations have variable lengths.

For agents, branching can include real operations, not only text. A planner can safely explore hypothetical patches in isolated workspaces, but cannot assume every external action is reversible. A search method should distinguish simulated states from committed environment states and prevent one branch's side effects from contaminating another. Replaying a tool call can be expensive or non-idempotent. The search graph should retain operation identity and observed outcomes; executor permission checks remain independent of controller confidence.

